\chapter{Performance considerations}\label{ch:perf}

\section{Where time is spent}
Model latency, not computation, dominates the duration of a request. A chat design makes one intake call and at least one narrative call, and up to three narrative calls when the number trace rejects a draft. The calculator and the verifier are pure computation. On the development machine used for this report (Python 3.14.8), 200 runs of \code{design_with_verification} on the section of Example~A in Appendix~\ref{app:example} took on average 0.23~ms per run, and ten calls of the 51-section explorer took on average 13.8~ms per call. These figures are single-machine measurements made for this report and are not guarantees.

\section{Time bounds}
\begin{itemize}
  \item Each model call times out at \code{LLM_TIMEOUT_S} (default 25~s). Each \code{generate()} call, across all models and retries, is limited to \code{LLM_BUDGET_S} (default 60~s). A request that makes several model calls can take longer than one budget.
  \item A model that times out is skipped, not retried; a model that returns 429, 500 or 503 gets one retry after a pause of at most 2 seconds.
  \item The Node proxy times out an upstream request after 300 seconds.
  \item The deep health check uses a budget of at most 20 seconds.
\end{itemize}

\section{Caching}
The live model check is cached for 60 seconds per model so that polling cannot spend quota. Frontend files are served with \code{Cache-Control: no-cache} and API responses with \code{no-store}. There is no result cache: identical designs are recomputed and logged again.

\section{Scalability notes}
\begin{itemize}
  \item The explorer evaluates up to 51 sections in one request with no model call.
  \item \code{GET /api/stats} uses \code{json_extract} over every row, so its cost grows linearly with the size of the log.
  \item The CSV export is capped at 5000 rows and the history list is paged (limit at most 200).
  \item The system is single process and single host (Chapter~\ref{ch:deploy}); a shared database would need a different \code{AuditLog} implementation.
\end{itemize}
